\documentclass[11pt]{article}
\usepackage[margin=1in]{geometry}
\usepackage{amsmath,amssymb,amsthm,mathtools}
\usepackage[T1]{fontenc}
\usepackage{lmodern,microtype}
\usepackage[colorlinks=true,allcolors=blue]{hyperref}
\newtheorem{theorem}{Theorem}
\newcommand{\R}{\mathbb R}
\newcommand{\ess}{\mathrm{ess}}
\newcommand{\dd}{\,\mathrm d}
\title{Robin half-space perturbations: a credited\\ dimension-dependent answer to AIM Problem 1.4}
\author{Alec Kriebel\\Independent researcher\\
\small ORCID: \href{https://orcid.org/0009-0001-9320-500X}{0009-0001-9320-500X}}
\date{Version 1.0 --- 14 October 2026}
\begin{document}
\maketitle
\begin{abstract}
We give a negative answer to the literal dimension-unrestricted wording of
AIM Problem 1.4 on attractive Robin bound states under compact half-space
perturbations. A nonzero smooth graph in $\R^4$, specified explicitly, has
spectrum $[-1,\infty)$ at Robin strength one. In contrast, every connected,
smooth, nontrivial compact half-space perturbation in ambient dimension two
or three binds below the flat threshold at every positive strength, including
overhangs and compact boundary components. Every compact Lipschitz graph in
dimension at least four has no below-threshold spectrum at sufficiently weak
strength. This is a credited exposition of standard mechanisms: Exner--Minakov
already proved planar binding and used the exponential/Stokes normal-defect
argument; the graph absence conclusion also follows from the established
Hoang--Hundertmark--Richter--Vugalter bound after the explicit geometric
comparison supplied here. We give both that implication and a self-contained
Hardy--trace certificate. Historical priority of the exact remaining geometric
formulations, novelty, worldwide open status and independent discovery remain
unresolved. Extensive AI assistance was used; this note is unrefereed.
\end{abstract}

\section{Question, conventions and attribution}
Write $n=m+1\ge2$, $z=(x,y)\in\R^m\times\R$, and $H=\{y>0\}$.
The attractive Robin operator $L_{\beta,\Omega}$, $\beta>0$, is the
self-adjoint operator associated with
\begin{equation}\label{form}
 q_{\beta,\Omega}[u]=\int_\Omega|\nabla u|^2\dd z
       -\beta\int_{\partial\Omega}|u|^2\dd S,\qquad D(q)=H^1(\Omega).
\end{equation}
Its outward condition is $\partial_\nu u=\beta u$; the flat threshold is
$-\beta^2$. The negative parameter $\alpha$ on the AIM source page
\cite{AIM} equals $-\beta$. The relevant problem and its immediately preceding
introduction impose no ambient-dimension restriction. Our negative answer is
to that literal wording, without a claim about an unstated intended dimension
or the question's current worldwide status.

A \emph{smooth compact perturbation} means a connected open $\Omega$ with
$C^\infty$ boundary, $\Omega\mathbin\triangle H$ bounded and $\Omega\ne H$.
We allow disconnected boundary and overhangs. For the graph statement we use
$f\in W^{1,\infty}_c(\R^m)$ and $\Omega_f=\{y>f(x)\}$.

\begin{theorem}\label{main}
The following statements hold.
\begin{enumerate}
\item Every smooth compact perturbation has
$\sigma_\ess(L_{\beta,\Omega})=[-\beta^2,\infty)$.
In ambient dimensions $n=2,3$, it has a discrete eigenvalue below
$-\beta^2$ for every $\beta>0$.
\item If $m\ge3$ and $f\in W^{1,\infty}_c(\R^m)$, there is
$\beta_0(f)>0$ such that, for $0<\beta\le\beta_0(f)$,
\[
 \sigma(L_{\beta,\Omega_f})=
 \sigma_\ess(L_{\beta,\Omega_f})=[-\beta^2,\infty).
\]
\item In $\R^4$, the nonzero $C^\infty$ graph $f=h/100$, where
\begin{equation}\label{bump}
 h(x)=\begin{cases}\exp(-1/(1-|x|^2)),&|x|<1,\\0,&|x|\ge1,\end{cases}
\end{equation}
satisfies $\sigma(L_{1,\Omega_f})=[-1,\infty)$.
For every prescribed $\beta_*>0$, a sufficiently small nonzero
$f=\varepsilon h$ gives another absence example at $\beta_*$.
\end{enumerate}
\end{theorem}

The planar single-infinite-$C^4$-curve, same-line-tail case of part 1 is due to
Exner--Minakov \cite[Theorem 5.2]{EM}. Their proof already uses
$\psi(x)e^{-\beta y}$ and Stokes' formula; its oriented curve factor is exactly
the planar specialization of our $-(1+\nu_y)$ defect. We credit the mechanism
as well as the theorem. The argument below verifies the additional boundary
components and dimension-three scope directly. These scopes' exact historical
priority remains unknown. Part 2 and the explicit example are also corollaries
of the published massive-relativistic bound \cite[Theorem A.6]{HHRV}, through
the all-form comparison in Section~\ref{prior}. We do not present the
dimensional threshold or the flat Robin/relativistic reduction as new
mechanisms. The latter is established in Frank--Laptev \cite{FL}; ground-state
null sequences and criticality are established frameworks \cite{PV}.

The theorem does not assert absence for one fixed nonzero graph at all
couplings, classify arbitrary nongraph perturbations in $n\ge4$, or exclude
embedded eigenvalues. Equality of spectral sets with the essential interval
alone does not exclude embedded point spectrum.

\section{Forms and essential spectrum}
Finite boundary charts on the altered compact part, together with the flat
trace estimate elsewhere, give for every $\epsilon>0$
\[
 \|\operatorname{Tr}u\|^2_{L^2(\partial\Omega)}
 \le \epsilon\|\nabla u\|^2+C_\epsilon\|u\|^2.
\]
Thus \eqref{form} is semibounded and closed, with a sufficiently shifted form
norm equivalent to $H^1$. Smooth functions compactly supported on the closure
are a form core in the smooth case. For Lipschitz graphs the bi-Lipschitz
flattening $(x,t)\mapsto(x,t+f(x))$ transfers density from $H$. Trial functions
need not satisfy the operator's Robin boundary condition.
Integration on flat vertical fibers gives
\begin{equation}\label{flat}
 Q_H[u]:=q_{\beta,H}[u]+\beta^2\|u\|^2
 =\|\nabla_xu\|^2+\|\partial_yu+\beta u\|^2\ge0.
\end{equation}
For any geometry in the theorem choose a smooth real partition
$\chi^2+\zeta^2=1$, with $\chi$ compactly supported and $\zeta$ zero near all
geometric changes. The function $\zeta u$, extended across the missing part
where necessary, lies in $H^1(H)$. IMS localization and semiboundedness of the
compact part imply, for a fixed bounded $K$,
\begin{equation}\label{ims}
 Q_\Omega[u]:=q_{\beta,\Omega}[u]+\beta^2\|u\|^2
 \ge -C\int_K|u|^2.
\end{equation}
Indeed the exterior form is nonnegative by \eqref{flat}; the compact form and
IMS error are bounded below by local $L^2$ mass.
If $\lambda<-\beta^2$ were essential, a normalized weakly null operator Weyl
sequence at $\lambda$ would have bounded $H^1$ norm by trace coercivity.
Local Rellich compactness makes its mass in $K$ tend to zero, contradicting
\eqref{ims} and $Q[u_j]\to\lambda+\beta^2<0$.

Conversely let $\eta\in C_c^\infty(\R^m)$ have norm one and choose
\[
 \eta_R(x)=R^{-m/2}\eta((x-a_R)/R),\qquad
 u_R(x,y)=\sqrt{2\beta}\,\eta_R(x)e^{ik\cdot x}e^{-\beta y}.
\]
Choose $a_R$ so the entire tangential support escapes to infinity and lies
outside the projection of every alteration. Extend by zero off those flat
fibers. Integration by parts against all $H^1$ tests shows $u_R$ is in the
operator domain: it has the flat Robin condition on its boundary support and
$L u_R=-\Delta u_R$. This uses no global $H^2$ domain characterization.
Its norm is one, it tends weakly to zero, and its residual at
$|k|^2-\beta^2$ has norm $O(|k|R^{-1})+O(R^{-2})$.
These operator Weyl sequences prove
\begin{equation}\label{ess}
 \sigma_\ess(L_{\beta,\Omega})=[-\beta^2,\infty).
\end{equation}
Below this interval spectrum is discrete, of finite multiplicity, and can
accumulate only at the threshold.

\section{Normal defect and low-dimensional binding}
Put $\phi=e^{-\beta y}$. Since $\Delta\phi=\beta^2\phi$ and
$\partial_\nu\phi=-\beta\nu_y\phi$, integration by parts gives, first for
compactly supported core functions $v$,
\begin{equation}\label{gs}
 Q_\Omega[\phi v]=\int_\Omega\phi^2|\nabla v|^2
 -\beta\int_{\partial\Omega}(1+\nu_y)\phi^2|v|^2\dd S.
\end{equation}
Real parts of the cross term give the same formula for complex $v$.
The integrand in
\[
 B_\beta(\Omega)=\int_{\partial\Omega}(1+\nu_y)e^{-2\beta y}\dd S
\]
is nonnegative and supported on the compact altered part. It is finite and
strictly positive. To see the latter, zero total would force
$\nu=-e_y$ everywhere by continuity. Each boundary component would then be
locally horizontal, open in its containing horizontal plane, and closed in
$\R^n$, hence the full plane. Compact agreement permits only $y=0$.
Connectedness and agreement with $H$ at infinity then force $\Omega=H$,
a contradiction. Overhangs and compact obstacle boundaries do not change
this sign or this argument.

Take $R_0$ larger than the projection of all alterations and a compactly
supported Lipschitz $\psi(x)$ equal to one near $B_{R_0}$.
The function $\phi\psi$ lies in $H^1(\Omega)$: the domain is bounded below,
and its upward tail is exponentially integrable. Vertical truncations above
the alteration have vanishing exponential errors, so \eqref{gs} applies.
The support of $\nabla\psi$ sees exactly flat whole vertical fibers, whence
\begin{equation}\label{capacity}
 Q_\Omega[\phi\psi]=\frac1{2\beta}\int_{\R^m}|\nabla\psi|^2
                         -\beta B_\beta(\Omega).
\end{equation}
For $m=1$ use two linear tapers of length $L$ outside $[-R_0,R_0]$;
their total energy is $2/L$. For $m=2$ use the radial logarithmic taper from
$R_0$ to $R$, of energy $2\pi/\log(R/R_0)$. Both tend to zero.
Equation~\eqref{capacity} is negative for a sufficiently large cutoff;
\eqref{ess} and min--max prove part 1.

This is an actual-domain null-sequence argument. In the terminology of
\cite{PV}, the auxiliary pair $P_0=-\Delta+\beta^2$,
$B_0u=\partial_\nu u+\beta\nu_yu$ has the positive solution $\phi$ and
transformed energy $\int_\Omega\phi^2|\nabla v|^2$.
The same tangential cutoffs, with vertical heights chosen afterwards to kill
each tail error and normalization on a fixed interior ball, supply its null
sequence. Smooth approximation supplies the required local core.
Subtracting the strictly positive compact boundary defect gives binding.
Thus we verify geometry and trace mass rather than assume flat criticality
transfers to an altered domain. The interior-potential perturbation remark in
\cite[Remark 5.22]{PV} is not a theorem about arbitrary domain changes.

\section{An elementary all-form graph estimate}\label{hardy}
Let $m\ge3$, $K=\operatorname{supp}f\subset\overline B_R$, and put
\[
 M=\|f\|_\infty,\quad G=\|\nabla f\|_\infty,\quad
 D=\sqrt{1+G^2}-1,\quad A=\frac{4R^2}{(m-2)^2},\quad
 c=\frac1{2(1+G^2)}.
\]
Classical tangential Hardy gives
\begin{equation}\label{localhardy}
 \int_K|a|^2\le A\int_{\R^m}|\nabla a|^2.
\end{equation}
For completeness, expand
$0\le\int|\nabla a+a_0p_\delta a|^2$ with
$a_0=(m-2)/2$ and $p_\delta=x/(|x|^2+\delta^2)$.
Since $\operatorname{div}p_\delta=((m-2)|x|^2+m\delta^2)/
(|x|^2+\delta^2)^2$, integration by parts yields
\[
 \int|\nabla a|^2\ge a_0^2\int
   \frac{|x|^2}{(|x|^2+\delta^2)^2}|a|^2.
\]
The omitted $a_0m\delta^2$ term is nonnegative for $m\ge3$.
Fatou and density prove Hardy; $1\le R^2/|x|^2$ on the ball gives
\eqref{localhardy}.

For a compact-support flattened core write
$u(x,t+f(x))=e^{-\beta(t+f(x))}w(x,t)$, $t>0$, and set
\[
 E_x=\int_H e^{-2\beta t}|\nabla_xw|^2,\quad
 E_t=\int_H e^{-2\beta t}|w_t|^2,\quad E=E_x+E_t.
\]
If $L_K=\int_{K\times(0,\infty)}e^{-2\beta t}|w|^2$
and $T_K=\int_K|w(x,0)|^2$, then
$L_K\le A E_x$ and vertical integration gives
\begin{equation}\label{trace}
 T_K=2\beta L_K-2\operatorname{Re}\int_{K\times(0,\infty)}
   e^{-2\beta t}\overline w w_t
 \le (2\beta A+\sqrt A)E.
\end{equation}
The last step uses $2\sqrt{E_xE_t}\le E_x+E_t$.
The weak chain rule and vertical-fiber integration give the exact identity
\begin{align}\label{graph}
 Q_f[u]={}&\int_H e^{-2\beta(t+f)}
    (|\nabla_xw-\nabla f\,w_t|^2+|w_t|^2)\nonumber\\
 &-\beta\int_{\R^m}e^{-2\beta f}
       (\sqrt{1+|\nabla f|^2}-1)|w(x,0)|^2\dd x.
\end{align}
In particular the Robin surface Jacobian is retained; the fiber endpoint
$+\beta|w(0)|^2$ combines with $-\beta\sqrt{1+|\nabla f|^2}|w(0)|^2$.
For $|g|\le G$, the triangle inequality squared gives
\[
 |a|^2+|b|^2\le2|a-gb|^2+(2G^2+1)|b|^2
 \le2(1+G^2)(|a-gb|^2+|b|^2).
\]
Equations~\eqref{trace}--\eqref{graph} therefore imply
\begin{equation}\label{lower}
 Q_f[u]\ge
 \big[e^{-2\beta M}c-\beta e^{2\beta M}D(2\beta A+\sqrt A)\big]E.
\end{equation}
Consequently
\begin{equation}\label{criterion}
 \beta D e^{4\beta M}(2\beta A+\sqrt A)\le c
\end{equation}
implies $Q_f\ge0$ on all of $H^1(\Omega_f)$. Indeed we first used a dense
compact-support core, where the exponential and its inverse are bounded,
then use form continuity and $H^1$ density. We never assume multiplication
by the globally growing $e^{\beta y}$ preserves $H^1$.
For fixed $f$, the left side of \eqref{criterion} tends to zero as
$\beta\downarrow0$; an interval $(0,\beta_0]$ follows. Together with
\eqref{ess} this proves part 2. The flat case $f=0$ holds for all strengths.
The homogeneous estimate \eqref{localhardy} is unavailable in $m=1,2$,
consistently with their capacity cutoffs.

\section{A complete earlier-theorem implication}\label{prior}
We now give the alternative derivation from the published
Hoang--Hundertmark--Richter--Vugalter theorem \cite{HHRV}, to make the
attribution substantive. Set
$U(x,t)=e^{\beta f(x)}u(x,t+f(x))$.
This is a bounded $H^1$ isomorphism because flattening is bi-Lipschitz and
$e^{\beta f}$ and its inverse are bounded Lipschitz multipliers.
Writing $b=U_t+\beta U$, direct expansion gives, on all of $H^1$ by density,
\begin{align*}
 Q_f[u]={}&\int_H e^{-2\beta f}
       (|\nabla_xU-\nabla f\,b|^2+|b|^2)\\
 &-\beta\int e^{-2\beta f}(\sqrt{1+|\nabla f|^2}-1)|U(x,0)|^2.
\end{align*}
Hence, with $a_\beta=e^{-2\beta M}c>0$,
$v_\beta=\beta D e^{4\beta M}/c$ and
$E_\beta[U]=\int_H(|\nabla_xU|^2+|U_t+\beta U|^2)$,
\begin{equation}\label{bridge}
 Q_f[u]\ge a_\beta\left(E_\beta[U]
                      -v_\beta\int_K|U(x,0)|^2\right).
\end{equation}
For the unitary tangential Fourier transform put
$\rho=(|p|^2+\beta^2)^{1/2}$ and $a_0=U(\cdot,0)$. Completion of the vertical
square gives
\begin{align}\label{square}
 E_\beta[U]&=\int_{\R^m}\int_0^\infty
       |\partial_t\widehat U+\rho\widehat U|^2\dd t\dd p
       +\int_{\R^m}(\rho-\beta)|\widehat a_0|^2\dd p\nonumber\\
 &\ge\langle a_0,T_\beta a_0\rangle,
 \qquad T_\beta=\sqrt{-\Delta+\beta^2}-\beta.
\end{align}
The identity extends from the core; $a_0\in H^{1/2}$.
Frank--Laptev \cite[Theorem 3.2 and Lemma 3.4]{FL} already establish the flat
Robin/relativistic duality, but their form-compact boundary-potential
hypothesis does not include the nondecaying attractive constant baseline.
We use \eqref{square} directly rather than apply that theorem outside its
hypotheses.

For fixed $\beta>0$ the form norm of $T_\beta+1$ is equivalent to
$H^{1/2}(\R^m)$. The bounded potential $V=-v_\beta\mathbf1_K$ is relatively
form small. Multiplication by $|V|^{1/2}$ is compact from $H^{1/2}$ to $L^2$,
by local Rellich compactness on a ball containing $K$. Thus the essential
spectrum stays $[0,\infty)$ and the hypotheses for the quantitative bound
in \cite{HHRV} are satisfied.

For $m\ge3$, published Theorem A.6 gives
\begin{equation}\label{hhrv}
 N(T_\beta+V)\le\frac{\theta^2}{(1-2\theta)^2}C_m
 \int(\theta^{-2}V_-)^{m/2}
      (\theta^{-2}V_-+2\beta)^{m/2},
\end{equation}
where $0<\theta<1/2$, $C_m=\frac m{m-2}\frac{\omega_m}{(2\pi)^m}$,
$\omega_m$ is the unit-ball volume, and $N$ counts strictly negative
eigenvalues with multiplicity. The version of record was published in 2022
and appears in the 2023 journal volume. Its source mass is our $\beta$ and
its dimension is our boundary dimension $m$. The 2017 arXiv v2
Theorem 3.6 instead has denominator $(1-4\theta)^2$ and
$0<\theta<1/4$. Our choice $\theta=1/4$ uses only the authenticated
\emph{published} formula; it must not be inserted in that preprint formula.

Take $\theta=1/4$. For $0<\beta\le1$ set $C_f=D e^{4M}/c$,
so $v_\beta\le\beta C_f$. Then
\begin{equation}\label{count}
 N(T_\beta-v_\beta\mathbf1_K)
 \le L_f\beta^m,\quad
 L_f=\frac{|K|C_m}{4}[16C_f(16C_f+2)]^{m/2}.
\end{equation}
For $D>0$ choose
$\beta_0=\min\{1,(2L_f)^{-1/m}\}$; for $D=0$, compact support makes
$f=0$. Since $N$ is an integer, \eqref{count} gives $N=0$.
There is no other negative spectrum, by the verified essential-spectrum
hypothesis; hence $T_\beta-v_\beta\mathbf1_K\ge0$.
Equations~\eqref{bridge}--\eqref{square} prove $Q_f\ge0$ on its entire
form domain. This is a complete earlier-theorem corollary for part 2,
not an effective-operator heuristic. Its exact geometric historical priority
remains unknown.

\section{Two exact certificates for the specified graph}
The flat extension of $e^{-1/s}$ at $s=0$ makes \eqref{bump} smooth.
It is positive at zero. For $r=|x|<1$ and $q=(1-r^2)^{-1}\ge1$,
\[
 |\nabla h|=2r q^2e^{-q}\le2q^2e^{-q}\le8/e^2<2.
\]
Thus for $f=h/100$, $M<1/100$, $G<1/50$, $D<1/5000$,
and $c>1250/2501$. At $\beta=1$ and $m\ge3$, $A\le4$ and
$2A+\sqrt A\le10$. The power-series comparison
$e^{1/25}<1/(1-1/25)=25/24$ gives
\[
 D e^{4M}(2A+\sqrt A)<1/480<1250/2501<c.
\]
The rational margin is $1250/2501-1/480=597499/1200480>0$.
This proves \eqref{criterion} exactly, without numerical spectral input.

Independently, for $m=3$, \eqref{hhrv} gives a certificate from the prior
bound. Here $v_1<2501/6000000$; putting $s=16v_1$ gives
$s<2501/375000<1/100$ and $s(s+2)<201/10000<1/40$.
Since $C_3=1/(2\pi^2)$ and $|K|\le4\pi/3$,
\[
 N(T_1-v_1\mathbf1_K)\le\frac{[s(s+2)]^{3/2}}{6\pi}<1/240<1.
\]
We used $s(s+2)<1$ and $\pi>1$ in the final coarse strict bound.
Both routes establish part 3 with \eqref{ess}. The graph and its domain
are connected and smoothly diffeomorphic to $\R^3$ and $H$, respectively;
$\Omega_f\mathbin\triangle H$ lies in a bounded unit cylinder.
For each fixed prescribed $\beta_*$, replacing $f$ by $\varepsilon h$ makes
$M,G\to0$ and $D=O(\varepsilon^2)$, so either certificate eventually applies.
The graph is allowed to depend on $\beta_*$.

\section{Source and verification limits}
The standard large-strength curvature binding theorem of
Pankrashkin--Popoff \cite{PP} is compatible with weak-strength absence.
The 2017 delta-interaction work and slides \cite{EKL} describe Robin
expectations, not the Robin theorems proved here. Deformation amplitude and
Robin strength are distinct parameters. Delta absence does not imply
selected-side Robin absence: the two-sided Robin form has the larger direct
sum form domain and its restriction to matching traces is the delta form
at strength $2\beta$. Its form order points the wrong way for that inference.
The full printed magnetic model of \cite{BSZ} requires $B>0$ comparable to
$\beta^\sigma$, $\sigma\ge2$, and sufficiently large $\beta$; it does not
apply to our zero-field statements.

The support kit preserves exact submitted and authored audit bytes, versioned
source-access records and finite standard-library checks. It excludes full
third-party bodies and image caches individually; access failures are not
evidence of novelty. The AIM passage was retrieved over plain HTTP without a
certificate bypass, a qualified provenance rather than authenticated HTTPS.
The published HHRV theorem was checked from the publisher body; the
Exner--Minakov and other specified arXiv interfaces are not represented as
byte-identical journal bodies. Priority of the exact remaining geometric
formulations, novelty, worldwide open status and independent discovery are
all unknown. No claim of first solution is made.

AI tools were used extensively in solving, checking, source comparison and
writing. This independent research note is unrefereed, has not received human
peer review, and is not a proof-assistant certificate. Its finite checks
verify selected exact algebra and custody; they do not replace the universal
arguments above. The accompanying edition records the required subsequent
whole-package adversarial review gate separately from historical math audits.
The paper and newly authored support materials are offered under CC BY 4.0.

\begin{thebibliography}{9}
\bibitem{AIM} American Institute of Mathematics, \emph{Shape optimization with
surface interactions}, Problem 1.4, \url{http://aimpl.org/shapesurface/1/}
(accessed 14 October 2026; literal dimension-unrestricted passage).
\bibitem{EM} P. Exner and A. Minakov, \emph{Curvature-induced bound states in
Robin waveguides and their asymptotical properties}, J. Math. Phys. 55 (2014),
122101, \href{https://doi.org/10.1063/1.4903184}{doi:10.1063/1.4903184}.
Inspected arXiv:1406.7624v1, Section 5, especially Theorem 5.2 and its proof.
\bibitem{HHRV} V. Hoang, D. Hundertmark, J. Richter and S. Vugalter,
\emph{Quantitative bounds versus existence of weakly coupled bound states
for Schr\"odinger type operators}, Ann. Henri Poincar\'e 24 (2023), 783--842,
\href{https://doi.org/10.1007/s00023-022-01228-3}{doi:10.1007/s00023-022-01228-3}.
Published online 10 September 2022; publisher Theorem A.6, printed p. 817.
The distinct arXiv:1610.09891v2 (2017) has different constants.
\bibitem{FL} R. L. Frank and A. Laptev, \emph{Spectral inequalities for
Schr\"odinger operators with surface potentials}, AMS Transl. Ser. 2,
225 (2008), 91--102. Inspected arXiv:0711.3473v1 and author-hosted version,
Theorem 3.2 and Lemma 3.4.
\bibitem{PV} Y. Pinchover and I. Versano, \emph{On criticality theory for
elliptic mixed boundary value problems in divergence form},
arXiv:2008.03699v1 (2020), especially Definition 3.15 and Theorem 6.4.
\bibitem{PP} K. Pankrashkin and N. Popoff, \emph{An effective Hamiltonian
for the eigenvalue asymptotics of the Robin Laplacian with a large parameter},
J. Math. Pures Appl. 106 (2016), 615--650,
\href{https://doi.org/10.1016/j.matpur.2016.03.005}{doi:10.1016/j.matpur.2016.03.005}.
Inspected arXiv:1502.00877v2, Corollary 1.6.
\bibitem{EKL} P. Exner, S. Kondej and V. Lotoreichik,
\emph{Asymptotics of the bound state induced by delta-interaction supported
on a weakly deformed plane}, arXiv:1703.10854v1 (2017), Section 4;
V. Lotoreichik, Trier 2017 slides, printed slides 9 and 15,
\url{https://aamp.fjfi.cvut.cz/lotoreichik/Trier2017.pdf}.
\bibitem{BSZ} D. Barseghyan, B. Schneider and Y. Zhang,
\emph{Spectral properties of the magnetic Robin Laplacian on curvilinear half-plane}, arXiv:2609.03783v1 (2026), Theorem 1.
The full printed positive-field and large-strength assumptions were checked
for applicability only, not every magnetic argument.
\end{thebibliography}
\end{document}
